\documentclass{article}

\usepackage[spanish]{babel}
\usepackage{amsmath,amssymb}

\title{Conceptos - Arquitectura de Computadores}
\author{Benjamín Rivas Beltrán}
\date{\today}

\begin{document}

\maketitle

\section{Conceptos básicos}

\begin{itemize}
  \item \textbf{Bit:} Unidad mínima de información que puede tomar el valor de 0 o 1 (encendido o apagado).
  \item \textbf{Byte:} Conjunto de 8 bits que representa un carácter o un valor numérico.
  \item \textbf{Información:} Todo aquello que puede reducuir la incertidumbre.
\end{itemize}

\section{Notas}

\begin{itemize}
  \item Los bytes son 8 bits, pues existen 256 combinaciones posibles de 0 y 1, lo que permite representar todos los caracteres de un alfabeto, incluyendo letras, números y símbolos.
\end{itemize}

\end{document}